%------------------------------------------------
\section{Заключение}

%------------------------------------------------
\begin{frame}{Заключение}
	\footnotesize % Уменьшаем общий шрифт
	\setlength{\leftmargini}{1.2em} % Сдвигаем буллиты влево, чтобы текст не переносился
	\begin{block}{1. Регулярная структура}
		\begin{itemize}
			\setlength{\itemsep}{0pt} % Убираем интервалы между пунктами
			\item $\tau_{IBS} \sim 1.5\times10^{4}$ с ($\approx 4$ ч) для $1.2 \times 10^{12}$ ppb.
			\item Стохастическое охлаждение \textbf{неэффективно}.
			\item Скачок критической энергии ($\gamma_{\text{tr}}$-jump) в барьерном ВЧ.
			\item Продольная микроволновая неустойчивость: снижение интенсивности в 2--4 раза.
			\item Необходима ВЧ-гимнастика выше критической энергии.
		\end{itemize}
	\end{block}
	
	\vspace{-0.5cm} % Подтягиваем нижние блоки выше
	\begin{columns}[t]
		\begin{column}{0.49\textwidth}
			\begin{block}{2. Резонансная структура}
				\scriptsize % Ещё чуть меньше для колонок
				\begin{itemize}
					\setlength{\itemsep}{0pt}
					\item $\tau_{\text{IBS}}$ в 3--4 раза ниже ($\approx 1$ ч) для $1.2 \times 10^{12}$ ppb.
					\item Без прохождения $\gamma_{\text{tr}}$.
					\item Модуляция градиентов квадруполей.
				\end{itemize}
			\end{block}
		\end{column}
		
		\begin{column}{0.49\textwidth}
			\begin{block}{3. Комбинированная структура}
				\scriptsize
				\begin{itemize}
					\setlength{\itemsep}{0pt}
					\item $\tau_{\text{IBS}}$ в 2 раза ниже ($\approx 2$ ч) для $1.2 \times 10^{12}$ ppb.
					\item СО \textbf{эффективно} при интенсивности $<5\times10^{10}$ ppb при энергии больше $6 \text{ ГэВ}$.  
					\item Без прохождения $\gamma_{\text{tr}}$
					\item Жесткая модуляция градиентов квадруполей.
					\item Оптимизация ДА -- секступольная коррекция.
				\end{itemize}
			\end{block}
		\end{column}
	\end{columns}
\end{frame}
%------------------------------------------------
\begin{frame}
	\frametitle{Благодарность}
	\begin{center}
		Спасибо за внимание!
	\end{center}
	
\end{frame}
